% analysis/pipeline_diagram-DE.tex
\section{Überblick über die AIM-Pipeline}

Dieser Abschnitt beantwortet eine Systemfrage: Welche Operatorfolge verbindet
CRS-Kontext, intrinsische Modellierung, Realisierung und Rückkopplung, ohne die
begriffliche Verantwortung zwischen den Ebenen zu verletzen?

Die Abbildung wird als Antwort durch Operatorreihenfolge eingeführt.

\begin{figure}[ht]
\centering
\begin{tikzpicture}[node distance=1.45cm,box/.style={draw,rectangle,rounded corners,minimum width=6.2cm,minimum height=0.95cm,align=center},arrow/.style={->,thick},feedback/.style={->,thick,dashed}]
\node[box] (crs) {CRS-Transformation \\ $\mathcal C$};
\node[box,below=of crs] (world) {World-Koordinaten};
\node[box,below=of world] (w2t) {World $\rightarrow$ Track \\ $\mathcal W$};
\node[box,below=of w2t] (model) {Krümmungsmodell \\ $\kappa(s)$ \\ (parametrisch + hybrid)};
\node[box,below=of model] (real) {Realisierung \\ $\mathcal R$};
\node[box,below=of real] (t2w) {Track $\rightarrow$ World \\ $\mathcal W^{-1}$};
\node[box,below=of t2w] (output) {Realisierte Geometrie};
\draw[arrow] (crs)--(world); \draw[arrow] (world)--(w2t);
\draw[arrow] (w2t)--(model); \draw[arrow] (model)--(real);
\draw[arrow] (real)--(t2w); \draw[arrow] (t2w)--(output);
\draw[feedback] (output.east)--++(1.2,0)|-node[pos=0.25,right]{Optimierung / SQP}(model.east);
\draw[feedback] (output.west)--++(-1.2,0)|-node[pos=0.25,left]{Messung / Daten}(w2t.west);
\end{tikzpicture}
\caption{AIM-Verarbeitungspipeline als Komposition von Operatoren, die Koordinatensysteme, Alignment-Modellierung, Realisierung und Optimierung verbinden.}
\end{figure}

Im Kapitelkontext begründet das Diagramm einen vorwärts gerichteten
Verarbeitungsstrang mit zwei ausdrücklichen Rückkopplungskanälen: Messdaten
stabilisieren die Projektionsinterpretation; Optimierungsrückkopplung
aktualisiert das Krümmungsmodell.

Daraus folgt, dass Pipeline-Korrektheit von der disziplinierten
Operatorkomposition und nicht von isolierten lokalen Berechnungen abhängt. Dies
bereitet die folgenden Abschnitte zu Interaktionsmustern und Fehlermodi vor.
